\documentclass{tietreport}

% Import images
\usepackage{graphicx}

% Bibliography configuration
\usepackage[
  date=year,
  style=alphabetic,
  backend=biber,
  sorting=nty,
  maxnames=5,
  minnames=3,
  maxalphanames=4,
  minalphanames=3,
  backref=true,
  doi=false,
  isbn=false,
  url=false,
  eprint=false
]{biblatex}
\addbibresource{references.bib}

%% ==================================================
%% PROJECT-SPECIFIC FIELDS -- CHARGEV
%% ==================================================

% Course/Lab header
\TitlePageHeader{%
  UCS503P - Software Engineering Project%
}

% Main project title
\ProjectTitle{%
  ChargEV: Find. Charge. Go Green.%
}

% Subtitle or project description
\ProjectSubTitle{%
  The Smart EV Charging Station Platform%
}

% Student information
\author{%
  \texttt{1024170103} Pehal \\
  \texttt{1024170431} Briti Arora%
}

% Group number, degree, year, program
\TitlePageSubText{%
  Group: BE Third Year, CSE%
}

% Faculty advisor/guide name
\AdvisorName{%
  Dr. Nisha Thakur%
}

% Evaluation period and department info
\TitlePageFooterText{{%
    \large Mid-Semester Evaluation \\[0.2cm]
    \bfseries Computer Science and Engineering
    Department%
  } \\[0.15cm]
  Thapar Institute of Engineering and Technology,
  Patiala%
}

%% ==================================================

\begin{document}

% Generate title page
\maketitle

% Table of contents (auto-generated from chapters)
\tableofcontents

% ===== CHAPTER 1: INTRODUCTION =====

\chapter{Introduction}

\section{Background and Context}

Electric vehicle (EV) adoption is accelerating, but the
charging infrastructure that supports it has not kept
pace in terms of usability. Drivers routinely do not
know which nearby stations are actually free, whether a
connector will be compatible with their vehicle, or
whether they will have to wait once they arrive. This
project, \textbf{ChargEV}, aims to remove this
uncertainty by giving EV drivers a single, reliable
platform to discover charging stations, check real-time
availability, and reserve a slot in advance -- turning
charging from a source of anxiety into a predictable,
routine task.

\subsection{Problem Statement}

While EV charging infrastructure continues to expand,
the tools available to EV drivers to actually use that
infrastructure remain fragmented and unreliable. The
specific problems this project addresses are:

\begin{itemize}
  \item \textbf{Discovery gap:} it is hard to find nearby
    charging stations from a single, reliable source.
  \item \textbf{Unclear availability:} station status
    (free, occupied, out of service) is rarely known
    before arriving.
  \item \textbf{Wasted time:} drivers wait at busy
    stations with no way to reserve a slot in advance.
  \item \textbf{Generic search results:} existing tools
    do not let drivers filter by connector type,
    distance, availability, and price together.
  \item \textbf{Difficult trip planning:} long-distance EV
    travel requires manually chaining together charging
    stops.
  \item \textbf{Low engagement:} there is little incentive
    for drivers to keep returning to a single platform.
\end{itemize}

As EV ownership grows, even small improvements in
charging predictability translate into measurable
reductions in range anxiety and wasted driver time.

\section{Project Objectives}

\begin{enumerate}
  \item \textbf{Centralized Station Locator:} Build a
    single source of truth for nearby EV charging
    stations, searchable by location.
  \item \textbf{Real-Time Availability \& Smart Booking:}
    Deliver a live availability dashboard and an online
    slot booking and reservation system so drivers can
    guarantee a charge before they leave.
  \item \textbf{Smart Search and Trip Planning:} Provide
    filters by connector type, distance, availability and
    price, and extend this into route-aware trip planning
    for long-distance travel.
  \item \textbf{Engagement Layer:} Implement a
    gamified rewards system (``EcoCoins'') that tracks
    savings and eco-impact to encourage repeat usage.
\end{enumerate}

% ===== CHAPTER 2: METHODOLOGY & DESIGN =====

\chapter{Methodology and Design}

\section{System Architecture}

\subsection{High-Level Design and Modelling}

Before writing any implementation code, the team invested
the first four weeks of the semester in requirement
analysis and system modelling, so that the frontend and
backend work started in parallel could stay consistent
with a shared blueprint. This involved:

\begin{itemize}
  \item A \textbf{UML Use Case Diagram} identifying the
    system boundary and its actors -- the EV Driver
    (primary actor), the Station Operator (secondary
    actor), and the Location Services API and MongoDB
    Atlas as supporting actors -- along with the main
    features: searching stations, applying smart filters,
    checking real-time availability, booking a slot,
    navigating to the station, planning multi-stop trips,
    and earning EcoCoins.
  \item An \textbf{Activity Diagram} mapping the
    step-by-step workflow of a station search, from
    entering a location and filters to displaying results
    on the map.
  \item A \textbf{Swimlane Activity Diagram} that
    distributes the complete search-book-charge flow
    across four responsibilities: the EV Driver/Frontend,
    the Backend API, Location Services (Maps/Places/
    Directions), and MongoDB Atlas -- covering everything
    from the initial search through booking, navigation,
    charging, and the final eco-impact summary.
  \item A \textbf{Sequence Diagram} describing the
    chronological, message-level interaction between the
    User, the Frontend, the Backend, the Database
    (MongoDB Atlas), and the Google APIs during a station
    search and booking.
  \item A \textbf{Collaboration Diagram} representing the
    same set of interactions as a structural map of
    communication links between those same components,
    with each link numbered to show the order of
    messages.
  \item A \textbf{Gantt Chart} laying out the full 20-week
    development plan, from requirement analysis through
    to the final report and demo.
\end{itemize}

All diagrams were cross-checked against the project
proposal to ensure consistency between the stated
requirements, the workflow, and the proposed system
architecture, and are presented in full in
Section~\ref{sec:diagrams}.

\subsection{Technical Stack}

\begin{itemize}
  \item \textbf{Frontend:} HTML5, CSS3, Bootstrap 5, and
    JavaScript, built for speed and cross-device
    responsiveness.
  \item \textbf{Backend:} Node.js and Express.js, chosen
    for high-performance, event-driven request routing.
  \item \textbf{Database:} MongoDB Atlas, a cloud-native
    document database.
  \item \textbf{Location Intelligence:} Google Maps API,
    Google Places API, and Google Directions API, used for
    station search, geocoding, and route planning.
\end{itemize}

\section{Implementation Details}

\subsection{Module 1: Frontend UI/UX and Prototype}

Frontend work began with research into interface
patterns used by other map-based and booking platforms,
which shaped the decision to use a map-plus-station-list
layout with a simplified booking flow of
\emph{station selection $\rightarrow$ slot selection
$\rightarrow$ confirmation}. Based on this research, a
fully navigable multi-page interface was implemented
using HTML5, CSS3 and Bootstrap 5, styled around a
verified, dark-themed EV charging brand identity:

\begin{itemize}
  \item A \textbf{Home page} with a hero search bar,
    quick filters (connector type, availability, 24/7
    open), and a live-style preview card for the nearest
    verified station.
  \item A dedicated \textbf{Stations page} combining a
    scrollable station directory with an interactive map
    view, so a driver can compare stations by price,
    connector type, and timings before choosing one.
  \item A \textbf{My Bookings} page listing a driver's
    scheduled charging sessions -- booking ID, station,
    connector, time window, duration, cost, and status --
    in a clean, tabular layout.
  \item A \textbf{Dashboard} presenting an active charging
    session with live-style telemetry (state of charge,
    charging rate, energy delivered, elapsed time, and
    accrued cost), alongside a saved-bookings table, in
    the format the platform will use once connected to
    the live backend.
  \item A \textbf{Create Account} flow for onboarding new
    drivers onto the Patiala EV charging network.
\end{itemize}

\begin{figure}[htbp]
  \centering
  \includegraphics[width=\textwidth]{images/chargev_home.jpeg}
  \caption{ChargEV Home Page: hero search, quick filters, and nearest-station preview.}
  \label{fig:ui-home}
\end{figure}

\begin{figure}[htbp]
  \centering
  \includegraphics[width=\textwidth]{images/chargev_stations_map.jpeg}
  \caption{Stations Page: Patiala EV charging directory with combined list and map view.}
  \label{fig:ui-stations}
\end{figure}

\begin{figure}[htbp]
  \centering
  \includegraphics[width=\textwidth]{images/chargev_my_bookings.jpeg}
  \caption{My Bookings: a driver's scheduled charging sessions in a clear tabular view.}
  \label{fig:ui-bookings}
\end{figure}

\begin{figure}[htbp]
  \centering
  \includegraphics[width=\textwidth]{images/chargev_dashboard.jpeg}
  \caption{Dashboard: live-style charging telemetry alongside saved booking records.}
  \label{fig:ui-dashboard}
\end{figure}

\begin{figure}[htbp]
  \centering
  \includegraphics[width=0.6\textwidth]{images/chargev_register.jpeg}
  \caption{Create Account: onboarding flow for new EV drivers.}
  \label{fig:ui-register}
\end{figure}

The layout was made \textbf{responsive} across desktop
and mobile screens using Bootstrap's grid system, and the
full navigation flow -- Home, Stations, Book Slot, My
Bookings, and Dashboard -- was validated end to end using
mock and simulated data, confirming that every screen
correctly renders the data structures the backend is
being built to supply.

\subsection{Module 2: Backend Development Setup}

In parallel, the initial backend environment was set up
so that frontend-backend integration can begin as soon as
the first API endpoints are ready:

\begin{itemize}
  \item The \textbf{Node.js and Express.js} server
    environment was initialised.
  \item The connection path between the \textbf{frontend,
    backend, and MongoDB Atlas} was explored and
    documented.
  \item The core backend modules were identified:
    \texttt{users}, \texttt{stations}, and
    \texttt{bookings}.
  \item Initial planning was done for the REST API
    structure that will expose these modules to the
    frontend in the next iteration.
\end{itemize}

At this stage, the backend is a scaffolded environment
rather than a working API -- no endpoints, database
schema, or live Google API calls have been implemented
yet. This is the next planned milestone (see
Section~\ref{sec:future-work}).

\section{System Modeling and Project Diagrams}
\label{sec:diagrams}

To provide a comprehensive blueprint of the ChargEV
system, the following diagrams were developed and
finalized during Weeks~2--4 of the semester.

\begin{figure}[htbp]
  \centering
  \includegraphics[width=\textwidth]{images/gantt_chart.pdf}
  \caption{ChargEV project timeline (Gantt Chart) across the 20-week semester.}
  \label{fig:gantt}
\end{figure}

\begin{figure}[htbp]
  \centering
  \includegraphics[width=\textwidth]{images/use_case_diagram.pdf}
  \caption{UML Use Case Diagram identifying the EV Driver, Station Operator, and supporting actors.}
  \label{fig:usecase}
\end{figure}

\begin{figure}[htbp]
  \centering
  \includegraphics[width=0.85\textwidth]{images/activity_diagram.pdf}
  \caption{UML Activity Diagram for the station search workflow.}
  \label{fig:activity}
\end{figure}

\begin{figure}[htbp]
  \centering
  \includegraphics[width=\textwidth]{images/activity_diagram_swimlanes.pdf}
  \caption{UML Activity Diagram mapped across swimlanes: EV Driver/Frontend, Backend API, Location Services, and MongoDB Atlas.}
  \label{fig:swimlanes}
\end{figure}

\begin{figure}[htbp]
  \centering
  \includegraphics[width=\textwidth]{images/sequence_diagram.pdf}
  \caption{UML Sequence Diagram depicting the chronological interaction between system components during a station search and booking.}
  \label{fig:sequence}
\end{figure}

\begin{figure}[htbp]
  \centering
  \includegraphics[width=0.9\textwidth]{images/collaboration_diagram.pdf}
  \caption{UML Collaboration Diagram highlighting the structural communication links between components.}
  \label{fig:collaboration}
\end{figure}

% ===== CHAPTER 3: RESULTS & EVALUATION =====

\chapter{Results and Evaluation}

\section{Testing Methodology}

As this is the mid-semester evaluation stage, our
efforts have successfully culminated in a validated
system blueprint and a fully navigable frontend
prototype for ChargEV:

\begin{itemize}
  \item \textbf{UML Design Validation:} The comprehensive
    UML models -- Use Case, Activity, Swimlane Activity,
    Sequence, and Collaboration diagrams -- have been
    cross-checked against the project proposal and
    validated as the logical blueprint for the
    discover-book-charge workflow.
  \item \textbf{UI/UX Validation:} The frontend has
    undergone rigorous interface validation across five
    core screens -- Home, Stations, My Bookings,
    Dashboard, and Create Account. The layout maintains
    responsiveness across desktop and mobile viewports,
    and input integrity is enforced through HTML5
    attributes and JavaScript event listeners.
  \item \textbf{Mock Integration \& Flow:} The interface
    transitions seamlessly from station discovery on the
    Home and Stations pages, through booking, to the
    Dashboard and My Bookings views. The mock and
    simulated data used at this stage successfully
    populates every station card, booking row, and
    telemetry widget, confirming that the frontend is
    fully capable of rendering the data structures the
    backend is being built to supply.
  \item \textbf{Backend Environment Validation:} The
    Node.js and Express.js server environment was
    initialised and validated as ready to host the
    \texttt{users}, \texttt{stations}, and
    \texttt{bookings} modules, with the MongoDB Atlas
    connection path confirmed ahead of full API
    development.
\end{itemize}

\section{Performance Metrics}

To keep development outcome-driven, the proposal defines
the metrics framework that will validate ChargEV's
booking workflow as backend integration proceeds:

\begin{table}[h]
\centering
\begin{tabular}{|l|p{6.5cm}|}
  \hline
  \textbf{Metric} & \textbf{Definition / Target} \\
  \hline
  Time-to-Charge (TTC) & Time from opening the app to
    starting a charge at a booked slot; target is a
    reduced median TTC versus a no-booking, walk-in
    baseline. \\
  \hline
  Booking Conversion Rate & Percentage of station
    searches that convert into a confirmed slot booking. \\
  \hline
  Availability Data Accuracy & Agreement between the
    displayed station status and the actual status
    observed on arrival. \\
  \hline
  User Engagement & Repeat-usage rate and EcoCoins
    accumulated per active user. \\
  \hline
\end{tabular}
\caption{Evaluation metrics framework for the ChargEV booking workflow.}
\label{tab:metrics}
\end{table}

This framework ensures that once the backend and live
station data are connected, ChargEV's success can be
measured objectively against the same discover-book-charge
workflow already validated in the UI and system diagrams.

\section{Analysis}

Our progress strictly adheres to our initial Gantt chart
objectives for the Requirements, UI/UX, and Backend Setup
phases shown in Figure~\ref{fig:gantt}. The robust,
multi-page, mock-data-driven design of the frontend
ensures that integrating the REST API will require
minimal UI refactoring, since every screen has already
been validated against the same data shapes the backend
modules are designed to return. The parallel system
modelling effort -- particularly the Swimlane Activity
Diagram and the Sequence and Collaboration diagrams --
gave both halves of the team a shared contract to build
against, so that the frontend and backend work completed
independently in the first four weeks will integrate
cleanly. Moving forward, the primary engineering focus
will shift to implementing the \texttt{stations} and
\texttt{bookings} REST endpoints and connecting them to
the already-validated frontend and location-intelligence
architecture.

% ===== CHAPTER 4: CONCLUSION =====

\chapter{Conclusion}

\section{Summary of Work}

To date, the team has:

\begin{itemize}
  \item Defined the problem scope and finalized the
    project proposal and technical stack.
  \item Completed a full set of UML and project
    diagrams -- Gantt Chart, Use Case, Activity,
    Swimlane Activity, Sequence, and Collaboration --
    validated against the proposal.
  \item Built and validated a fully navigable, responsive
    frontend prototype spanning the Home, Stations, My
    Bookings, Dashboard, and Create Account screens.
  \item Initialised and validated the Node.js/Express
    backend environment and its planned connection to
    MongoDB Atlas.
\end{itemize}

\section{Key Findings}

\begin{itemize}
  \item \textbf{Finding 1:} Building the frontend first,
    in the shape defined by the Use Case and Sequence
    diagrams, means the planned REST API can be dropped
    in with minimal UI refactoring later.
  \item \textbf{Finding 2:} The Swimlane Activity Diagram
    clarified responsibility boundaries between the
    Backend API and Location Services early, before any
    backend code was written, which directly shaped how
    the \texttt{stations} and \texttt{bookings} modules
    were scoped.
  \item \textbf{Finding 3:} Designing the Dashboard and My
    Bookings screens around the exact fields the MongoDB
    schema will store (booking ID, connector, time window,
    cost, status) gives the team a validated target schema
    to implement against in the next iteration.
\end{itemize}

\section{Future Work and Recommendations}
\label{sec:future-work}

Following the established Gantt chart timeline, the
immediate next steps are:

\begin{enumerate}
  \item \textbf{REST API Development:} Implement the
    \texttt{stations}, \texttt{bookings}, and
    \texttt{users} endpoints in Express and connect them
    to a MongoDB Atlas schema.
  \item \textbf{Location Intelligence Integration:}
    Integrate the Google Maps, Places, and Directions
    APIs for live station search, geocoding, and routing.
  \item \textbf{Frontend-Backend Integration:} Connect the
    validated frontend screens to live API responses, so
    the already-built UI is driven end-to-end.
  \item \textbf{Trip Planning and EcoCoins:} Build the
    route-aware trip planning module and the EcoCoins
    rewards and eco-impact tracking engine.
  \item \textbf{Testing and CI/CD:} Add unit and
    integration tests, and set up a CI/CD pipeline for
    staged, incremental releases.
\end{enumerate}

\section{Final Remarks}

The first four weeks of this project have established a
consistent, validated blueprint -- across the proposal,
the full UML diagram set, and a fully navigable frontend
prototype -- for the discover-book-charge workflow that
ChargEV is built around. This foundation is intended to
make the coming integration phase, where the frontend,
backend, and external APIs are connected end-to-end,
significantly more predictable.

% ===== BIBLIOGRAPHY =====

% Print bibliography with heading in TOC
\printbibliography[heading=bibintoc]

\end{document}
